\chapter{Introducción}

En ensayos repetitivos o de alta densidad de datos, la operación manual del instrumental introduce
fatiga en los operadores, tiempos muertos prolongados y una dispersión estadística apreciable atribuible a demoras de
lectura y transcripción. Frente a estas limitaciones, los sistemas de medición automatizados permiten gobernar los
parámetros operativos de generadores de señal y sistemas de adquisición mediante software, garantizando secuencias de
prueba reproducibles, sincronización determinística y captura masiva de datos en tiempo real.

En este trabajo práctico se aborda el diseño e implementación de un sistema automatizado de control instrumental basado
en una computadora, interfaz RS-232 y el protocolo estandarizado SCPI (\emph{Standard Commands for Programmable
Instruments}), orientado a la caracterización experimental de la respuesta en frecuencia y el factor de calidad ($Q$) de
un circuito resonante serie RLC.

\section{Objetivos del Trabajo Práctico}
\begin{itemize}
    \item Diseñar e implementar una aplicación para el control remoto de un generador de funciones
    comercial a través de su interfaz serie RS-232 mediante comandos SCPI.
    \item Desarrollar algoritmos de barrido en frecuencia discretos, contemplando modalidades de avance automático
    con temporización configurable y avance manual con validación por teclado.
    \item Caracterizar experimentalmente la respuesta en frecuencia de un circuito resonante RLC serie provisto por la
    cátedra, identificando la frecuencia de resonancia ($f_0$) y las frecuencias de corte a $\qty{-3}{\decibel}$.
    \item Cuantificar el factor de calidad o sobretensión ($Q$) del sistema.
    \item Validar la curva de respuesta en frecuencia en el dominio de frecuencia, evaluando el impacto del ventaneo y
    el modo de persistencia infinita en pantalla.
\end{itemize}

\section{Fundamentación Teórica}

\subsection{Arquitectura de Instrumentación Programable y Protocolo SCPI}
Históricamente, la interconexión de instrumentos de prueba surgió con el bus HPIB (\emph{Hewlett-Packard Interface
Bus}), estandarizado posteriormente como IEEE-488.1 (GPIB). Dicha norma definió las especificaciones mecánicas,
eléctricas y de protocolo básico de enlace paralelo. A fin de unificar la sintaxis y los códigos de operación entre
diferentes fabricantes, se promulgaron el estándar IEEE-488.2 y la especificación SCPI (\emph{Standard Commands for
Programmable Instruments}).

La norma SCPI define una estructura de comandos basada en cadenas ASCII organizadas de forma jerárquica como un árbol.
El árbol posee un nodo raíz, del cual se desprenden trayectorias hacia nodos secundarios y nodos
hoja. La separación entre niveles jerárquicos se efectúa mediante dos puntos (\texttt{:}), y los
parámetros asociados a cada instrucción se separan del comando por un espacio en blanco. Los comandos aceptan una
forma extendida y una forma corta abreviada (resaltada en mayúsculas en los manuales de referencia):
\begin{itemize}
    \item \emph{Comandos de configuración}: Modifican un estado interno o parámetro del instrumento, por ejemplo
    \texttt{:FREQ 10000} o \texttt{:AMPL:VOLT 2.5}.
\item \emph{Comandos de consulta}: Solicitan al equipo el estado o valor de un parámetro,
    finalizando siempre con el carácter de interrogación \texttt{?}, como \texttt{:FREQ?} o \texttt{*IDN?}.
    \item \emph{Comandos comunes IEEE-488.2}: Comienzan con un asterisco y realizan acciones globales
    predefinidas: \texttt{*RST} (reinicio a condiciones de fábrica), \texttt{*IDN?} (consulta de fabricante, modelo,
    número de serie y versión de firmware) y \texttt{*CLS} (borrado de registros de estado).
\end{itemize}

\subsection{Teoría del Circuito Resonante Serie RLC}
Un circuito resonante serie está compuesto por una resistencia total equivalente $R_{\text{tot}}$, una inductancia
concentrada $L$ y una capacitancia $C$, excitados por una fuente senoidal de tensión $v_s(t) = V_s \cos(\omega t)$.
La impedancia compleja equivalente del conjunto vista desde los bornes de la fuente se expresa como:
\begin{equation}
    Z(j\omega) = R_{\text{tot}} + j \left( \omega L - \frac{1}{\omega C} \right)
    \label{eq.ej1:impedancia}
\end{equation}
donde $R_{\text{tot}}$ engloba tanto las resistencias explícitas del circuito como la resistencia interna del
generador ($R_g$) y las pérdidas óhmicas parásitas del devanado de la bobina ($R_L$).

La condición de resonancia serie se produce en $\omega_0$ en la cual la reactancia inductiva
$X_L = \omega L$ cancela con exactitud a la reactancia capacitiva $X_C = \frac{1}{\omega C}$, anulando la parte
imaginaria de la impedancia:
\begin{equation}
    \omega_0 L - \frac{1}{\omega_0 C} = 0 \implies \omega_0 = \frac{1}{\sqrt{LC}}
    \label{eq.ej1:w0}
\end{equation}
Por consiguiente, la frecuencia de resonancia teórica en Hertz queda determinada por la relación de Thomson:
\begin{equation}
    f_0 = \frac{1}{2 \pi \sqrt{L C}}
    \label{eq.ej1:fres}
\end{equation}

A la frecuencia $f_0$, el circuito manifiesta las siguientes propiedades singulares:
\begin{enumerate}
    \item La impedancia total alcanza su valor mínimo absoluto y puramente resistivo: $Z(j\omega_0) = R_{\text{tot}}$.
    \item La corriente que circula por la malla adquiere su módulo máximo: $I_0 = \frac{V_s}{R_{\text{tot}}}$,
    encontrándose en fase con la tensión de excitación ($\Delta \varphi = 0^\circ$).
    \item Las caídas de tensión en los componentes reactivos individuales ($v_L$ y $v_C$) presentan una magnitud que
    puede superar ampliamente el valor de la tensión de entrada $V_s$. Este fenómeno se denomina sobretensión o
    resonancia de tensión.
\end{enumerate}

\subsection{Factor de Calidad y Relación con el Ancho de Banda}
El factor de calidad $Q$ (o factor de selectividad) cuantifica la relación entre la energía reactiva almacenada en el
sistema y la energía activa disipada por ciclo en los elementos resistivos. Para una malla resonante serie, se define
como:
\begin{equation}
    Q = \frac{\omega_0 L}{R_{\text{tot}}} = \frac{1}{\omega_0 C R_{\text{tot}}} = \frac{1}{R_{\text{tot}}} \sqrt{\frac{L}{C}}
    \label{eq.ej1:q_definicion}
\end{equation}

Analizando la caída de tensión en bornes del capacitor a la frecuencia de resonancia:
\begin{equation}
    V_C(\omega_0) = I_0 \cdot X_C(\omega_0) = \left( \frac{V_s}{R_{\text{tot}}} \right) \left( \frac{1}{\omega_0 C} \right) = Q \cdot V_s
    \label{eq.ej1:magnificacion}
\end{equation}
La ecuación~\eqref{eq.ej1:magnificacion} evidencia que la relación entre la tensión pico a pico en el capacitor y la
tensión total aplicada a la rama serie en resonancia proporciona una estimación directa del factor de calidad del
sistema, denominado factor de sobretensión:
\begin{equation}
    Q_v = \frac{V_C(f_0)}{V_{\text{in}}(f_0)}
    \label{eq.ej1:qv}
\end{equation}

Por otra parte, si se analiza la curva de transferencia en función de la frecuencia, se definen las frecuencias de
corte inferior ($f_1$ o $f_{\text{inf}}$) y superior ($f_2$ o $f_{\text{sup}}$) como aquellos puntos donde la potencia
entregada cae a la mitad de su valor de resonancia ($\qty{-3}{\decibel}$ respecto al máximo). En dichos puntos, la
tensión en el capacitor se reduce a un factor $1/\sqrt{2} \approx 0.7071$ del valor pico:
\begin{equation}
    V_C(f_1) = V_C(f_2) = \frac{V_C(f_0)}{\sqrt{2}}
    \label{eq.ej1:umbral_3db}
\end{equation}

El ancho de banda a $\qty{-3}{\decibel}$ se define como la diferencia entre ambas frecuencias de corte:
\begin{equation}
    \Delta f = f_{\text{sup}} - f_{\text{inf}}
    \label{eq.ej1:bw}
\end{equation}
El factor de calidad experimental derivado de la selectividad en frecuencia resulta entonces:
\begin{equation}
    Q_{\Delta f} = \frac{f_0}{\Delta f}
    \label{eq.ej1:q_bw}
\end{equation}

